\section{Artifact interfaces and sealed evaluation}
\label{app:artifact-interfaces}

The artifact exposes three versioned interfaces:

\begin{itemize}[leftmargin=1.5em,itemsep=0.1em,topsep=0.2em]
  \item \path{observerbench.api.v0} for control composition;
  \item \path{observerbench.effect_prediction.v0} for held-out finite-effect
  prediction; and
  \item \path{observerbench.safety_protocol.v0} for safety triage.
\end{itemize}

Their machine-readable schemas, example submissions, and validation commands
ship with the repository. The version identifiers prevent a result produced
under one contract from being compared silently with a result produced under
another.

The task registry states which submission route each bundled task supports.
In the current release, the analytic control task provides the checked
external-observer adapter. Its benchmark wrapper constructs the controller and
reuses the task's fixed generator, readouts, normalization, and metrics. A
lower-level run may replace the controller, but the resulting score is marked
non-comparable because it no longer represents the same decision problem.

Public safety-task packs contain the threat model, allowed actions,
consequence loss, operating budgets, training measurements, and held-out
queries. They do not contain the held-out evaluator answers. The evaluator
stores those targets separately and links them to the public pack by a
cryptographic hash. A submission can therefore be checked against the declared
task while the labels used for final scoring remain sealed. This layout also
allows a researcher to add a data-only safety task without modifying the core
workbench.

The bundled GPT-2 IOI and Qwen induction-copy registries implement the same
separation for finite-effect prediction. Each registry exposes frozen
measurement tables at declared budgets, checks the model and split recorded by
the task, and uses different prompts for calibration measurements and held-out
targets. An outside observer can submit its predictions without loading the
model or rerunning the intervention experiment. The registries include
per-component and additive ridge observers as initial comparison rows.

This is a data-only boundary, not a requirement to execute contributor code.
For effect prediction, a submission supplies one row per held-out intervention
with the fields \path{schema_version}, \path{query_id}, and
\path{predicted_effect}, together with an ObserverCard that states the method
and its access assumptions. The evaluator rejects missing, duplicate,
unexpected, or non-finite predictions, joins accepted rows to sealed targets,
and returns aggregate prediction metrics. It does not yet select an action or
return action loss for an arbitrary submitted table. The action-selection
studies in this paper use the frozen IOI and Qwen experiment runners, which
fit the mean-effect or direct-loss observer, apply the fixed candidate-pool
selector, and score the selected masks on held-out prompts. Their action
results should not be read as action results for the submitted AtP rows.
The next extension must join these paths while preserving the declared
estimand, target, candidate pools, no-op, and tie rule. It must return the
selected action, action loss, and comparisons against no-op and named
reference observers, not infer decision quality from MAE. Keeping submissions inert
and targets private protects scoring labels and avoids running outside code.
It does not make the underlying telemetry trustworthy, test adaptive evasion,
or prevent repeated tuning to a public panel.

The external APPS registry contains separate Qwen2.5-7B, Gemma-2-9B-it, and
Qwen3.5-9B scorecards. Prompt, output, text, dense-residual, and SAE-feature
observers are evaluated under the same audit contract within each model panel.
The Qwen3.5 release includes the frozen source manifest and an aggregate
checked summary containing the scorecard rows, paired contrasts, and official
SAE selection. Per-example evaluator labels remain sealed. The registry does
not issue a cross-model rank.

Runs use strict JSON and record
the contract and result-schema versions, component classes and parameters,
package and dependency versions, the source revision, configuration and result
hashes, and observer provenance supplied by the submitter.
